\section{Нормальный (гауссовский) случайный вектор.}
\begin{definition}
    Пусть есть случайный вектор $\vec{\xi}$, с математическим ожиданием $\vec{m} = \Ex \vec{\xi}$. \\
    Рассмотрим матрицу ковариации этого случайного вектора, которая определяется как $R = \{\cov \xi_i, \xi_j\}_{i, j = 1, 1}^{n, n}$.
    Будем говорить, что $\vec{\xi}$ -- нормальный случайный вектор с параметрами $\vec{m}$ и $R$, если $\det R \neq 0$ и он имеет плотность:
    \[
        p_{\vec{\xi}}(\vec{x}) = \dfrac{1}{(\sqrt {2\pi})^n \sqrt {|\det R|}} \exp\biggr( -\dfrac{1}{2}(\vec{x} - \vec{m})^T R^{-1}(\vec{x} - \vec{m}) \biggr).
    \]
\end{definition}
\begin{lemma}
    Заметим, что для нормального случайного вектора некоррелированность и независимость компонент эквивалентна -- совместная плотность просто представляется в виде произведения плотностей.
\end{lemma}
\begin{note}
    Заметим, что если $\xi \sim N(\vec{m}, R)$ и $\eta = B \xi$, то плотность преобразуется таким образом, что $\eta \sim N(B \vec{m}, BRB^T)$:
    Это так, поскольку $\xi = B^{-1}\eta$ и \[|J| = \dfrac{1}{|\det B|} = \dfrac{1}{\sqrt {\det B \det B^T}}.\]
    Теперь прямо посмотрим на плотность $\eta$:
    \[
        p_{\vec{\eta}}(\vec{x}) = \dfrac{1}{(\sqrt{2\pi})^{n} \sqrt{\det|BRB^T|}} \exp (-\dfrac{1}{2}(B^{-1}\vec{x} - \vec{m})^TR^{-1}(B^{-1}\vec{x} - \vec{m})).
    \]
    Выражение в показателе экспоненты преобразуется как
    \[
        (B^{-1}(\vec{x} - B\vec{m}))^T R^{-1} B^{-1}(\vec{x} - B\vec{m}) = (\vec{x} - B\vec{m})(BRB^T)^{-1}(\vec{x} - B\vec{m}).
    \]
    что и показывает справедливость нашего утверждения.
\end{note}
\subsection{Характеристическая функция многомерного гауссовского распределения.}
\begin{definition}
    Модифицируем определение характеристической функции на случайные вектора.
    Собственно, для этого не надо выдумывать что-то сверхестественное -- просто используем многомерное преобразование Фурье:
    \[
        \phi_{\vec{\xi}}(\vec{t}) = \Ex e^{i \langle \vec{t}, \vec{\xi} \rangle}.
    \]
\end{definition}
\begin{lemma}
    Характеристическая функция невырожденного гауссовского случайного вектора $\vec{\xi} \sim N(\vec{m}, R)$ имеет вид
    \[
        \phi_{\vec{\xi}}(\vec{t}) = e^{i \langle \vec{m}, \vec{t} \rangle - \frac{1}{2} \vec{t}^T R \vec{t}}.
    \]
\end{lemma}
\begin{proof}
    В силу положительной определённости матрицы $R$ существует такая ортогональная $U$, что $URU^T = diag(\sigma_1^2, \ldots \sigma_n^2)$.
    Рассмотрим тогда случайный вектор $\vec{\eta} = U\vec{\xi}$.
    Он распределён как $N(U\vec{m}, diag(\sigma_1^2, \ldots, \sigma_n^2))$.
    Его компоненты -- некоррелированные случайные величины, а следовательно -- независимые.
    Тогда его харфункция является произведение харфункций его компонент и равна
    \[
        \phi_{\vec{\eta}}(\vec{t}) = e^{i \langle U\vec{m}, \vec{t} \rangle- \frac{1}{2} \vec{t}^T diag(\sigma_1^2, \ldots) \vec{t}}.
    \]
    Поскольку мы знаем как меняются харфункции при линейных преобразованиях случайных величин (а, следовательно, и случайных векторов -- это несложно показать по определению), обратное преобразование $\vec{\xi} = U^{-1}\vec{\eta}$ даст нам утверждение леммы.
\end{proof}
\begin{note}
    Характеристические функции случайных векторов во многом имеют свойства, аналогичные свойствам характеристических функций случайных величин.
\end{note}
\subsection{Обобщённый нормальный вектор. Линейное преобразование обобщённого нормального вектора.}
\begin{definition}
    Требование невырожденности матрицы ковариации может быть слегка излишним, поэтому мы можем сделать обобщение, и называть $\vec{\xi}$ обобщённым нормальным вектором, если
    \[
        \phi_{\vec{\xi}}(\vec{t}) = e^{i \langle \vec{t}, \vec{m} \rangle - \frac{1}{2}\vec{t}^T R \vec{t}}.
    \]
    для некоторого вектора $\vec{m}$ и неотрицательно определённой симметрической матрицы $R$.
\end{definition}
\begin{lemma}
    Пусть $\vec{\xi} \sim (\vec{m}, R)$ -- обобщённый гауссовский вектор. Тогда если $\vec{\eta} = A\vec{\xi} + \vec{b}$, то
    $\vec{\eta}$ -- также обобщённый нормальный вектор, который распределён $\sim N(A\vec{m} + \vec{b}, ARA^T)$.
\end{lemma}
\begin{proof}
    \[
        \phi_{\vec{\eta}}(\vec{t}) = \Ex e^{i \langle \vec{t}, A\vec{\xi} + \vec{b} \rangle} = e^{i\langle \vec{t}, \vec{b} \rangle} \Ex e^{i \vec{t}^T A \vec{\xi}} = e^{i\langle \vec{t}, \vec{b} \rangle} \Ex e^{i (A^T\vec{t})^T \vec{\xi}}
        = e^{i \langle \vec{t}, \vec{b} \rangle} \phi_{\vec{\xi}}(A^T \vec{t}) = e^{i \langle \vec{t}, \vec{b} \rangle} \cdot e^{i \langle \vec{t}, A\vec{m} \rangle} \cdot e^{-\frac{1}{2} \vec{t}^T ARA^T \vec{t}}.
    \]
\end{proof}